\documentclass[12pt]{article}
\usepackage[utf8]{inputenc}
\usepackage[T1]{fontenc}
\usepackage[margin=1in]{geometry}
\usepackage{setspace}
\usepackage{xcolor}
\usepackage[colorlinks=true, urlcolor=blue!60!black, linkcolor=blue!60!black]{hyperref}
\setlength{\parindent}{2em}
\setlength{\parskip}{1.25em}
\renewcommand{\baselinestretch}{1.0}
\begin{document}
\begin{center}
{\LARGE \textit{Printz v. United States}, 521 U.S. 898 (1997)\par}
\vspace{0.8em}
{\large Case Brief\par}
\end{center}
\vspace{1em}

\section{Memory Jogger}
In \href{https://www.courtlistener.com/opinion/118148/printz-v-united-states/}{\textit{Printz v. United States}}, the Supreme Court struck down the Brady Act's interim mandate requiring local law enforcement officers to conduct handgun background checks, establishing a categorical anti-commandeering rule that Congress may not conscript state or local executive officers to administer or enforce a federal regulatory program.

\section{Detailed Case Facts}
In 1993, Congress amended the Gun Control Act of 1968, 18 U.S.C. \S\ 921 \textit{et seq.}, by enacting the Brady Handgun Violence Prevention Act, Pub. L. No. 103-159, 107 Stat. 1536 (\href{https://www.law.cornell.edu/supct/html/95-1478.ZO.html}{Cornell LII, Majority Opinion}). The Act directed the U.S. Attorney General to establish a national instant criminal background check system (NICS) by November 30, 1998, and instituted mandatory interim procedures to govern handgun sales until that federal system became operational (\href{https://www.law.cornell.edu/supct/html/95-1478.ZO.html}{Cornell LII, Majority Opinion}).

Under the interim provisions, a federally licensed firearms dealer proposing to transfer a handgun was required to obtain a sworn statement from the prospective buyer (the ``Brady Form''), verify the buyer's identity, forward notice of the contents and a copy of the Brady Form to the ``chief law enforcement officer'' (CLEO) of the buyer's residence under 18 U.S.C. \S\S\ 922(s)(1)(A)(i)(III)--(IV), and wait five business days before completing the sale unless the CLEO approved it earlier under \S\ 922(s)(1)(A)(ii) (\href{https://www.law.cornell.edu/supct/html/95-1478.ZO.html}{Cornell LII, Majority Opinion}). In states without qualifying state permit or instant-check alternatives, \S\ 922(s)(2) commanded the local CLEO to ``make a reasonable effort to ascertain within 5 business days whether receipt or possession would be in violation of the law, including research in whatever State and local record keeping systems are available and in a national system designated by the Attorney General'' (\href{https://www.law.cornell.edu/supct/html/95-1478.ZO.html}{Cornell LII, Majority Opinion}). If a CLEO found no basis to object, \S\ 922(s)(6)(B)(i) required the CLEO to destroy the Brady Form and related records; if the CLEO notified a dealer that a buyer was ineligible, \S\ 922(s)(6)(C) required the CLEO, upon request, to provide the buyer a written statement of reasons (\href{https://www.law.cornell.edu/supct/html/95-1478.ZO.html}{Cornell LII, Majority Opinion}).

Two county sheriffs serving as CLEOs---Jay Printz, Sheriff/Coroner of Ravalli County, Montana, and Richard Mack, Sheriff of Graham County, Arizona---objected to being pressed into federal service without state authorization and filed separate federal lawsuits challenging the constitutionality of the Brady Act's interim CLEO mandates (\href{https://www.law.cornell.edu/supct/html/95-1478.ZO.html}{Cornell LII, Majority Opinion}; \href{https://www.oyez.org/cases/1996/95-1478}{Oyez}).

\section{Procedural History}
Sheriff Printz filed suit in the U.S. District Court for the District of Montana, and Sheriff Mack filed a separate action in the U.S. District Court for the District of Arizona (\href{https://www.law.cornell.edu/supct/html/95-1478.ZO.html}{Cornell LII, Majority Opinion}). Both district courts held that the interim provision compelling CLEOs to perform background checks violated the Tenth Amendment, but concluded that the background-check mandate was severable from the remainder of the Brady Act, effectively leaving a voluntary background-check system in place: \textit{Printz v. United States}, 854 F. Supp. 1503 (D. Mont. 1994), and \textit{Mack v. United States}, 856 F. Supp. 1372 (D. Ariz. 1994) (\href{https://www.law.cornell.edu/supct/html/95-1478.ZS.html}{Cornell LII, Syllabus}; \href{https://www.law.cornell.edu/supct/html/95-1478.ZO.html}{Cornell LII, Majority Opinion}).

On consolidated appeal, a divided panel of the U.S. Court of Appeals for the Ninth Circuit reversed, holding that none of the Brady Act's interim provisions were unconstitutional, 66 F.3d 1025 (9th Cir. 1995), with Judge Fernandez dissenting (\href{https://www.law.cornell.edu/supct/html/95-1478.ZO.html}{Cornell LII, Majority Opinion}). The U.S. Supreme Court granted certiorari in both cases---No. 95-1478 (\textit{Printz v. United States}) and No. 95-1503 (\textit{Mack v. United States})---heard oral argument on December 3, 1996, and issued its decision on June 27, 1997 (\href{https://www.law.cornell.edu/supct/html/95-1478.ZS.html}{Cornell LII, Syllabus}; \href{https://www.oyez.org/cases/1996/95-1478}{Oyez}).

\section{Judicial Votes}
The Supreme Court reversed the Ninth Circuit by a \textbf{5--4} vote (\href{https://www.law.cornell.edu/supct/html/95-1478.ZS.html}{Cornell LII, Syllabus}; \href{https://www.courtlistener.com/opinion/118148/printz-v-united-states/}{CourtListener}):
\begin{itemize}
\item \textbf{Majority (5):} Justice Antonin Scalia delivered the opinion of the Court, joined in full by Chief Justice William H. Rehnquist, Justice Sandra Day O'Connor, Justice Anthony M. Kennedy, and Justice Clarence Thomas.
\item \textbf{Concurring Opinions (2):} Justice Sandra Day O'Connor filed a concurring opinion; Justice Clarence Thomas filed a concurring opinion.
\item \textbf{Dissenting Opinions (4 Justices, 3 Opinions):} Justice John Paul Stevens filed the principal dissenting opinion, joined by Justice David H. Souter, Justice Ruth Bader Ginsburg, and Justice Stephen G. Breyer. Justice David H. Souter filed a separate dissenting opinion. Justice Stephen G. Breyer filed a separate dissenting opinion, joined by Justice John Paul Stevens.
\end{itemize}

\section{Holding}
\textbf{Judgment:} Reversed (\href{https://www.law.cornell.edu/supct/html/95-1478.ZS.html}{Cornell LII, Syllabus}).

\textbf{Rule Announced:} The Federal Government may neither issue directives requiring the states to address particular problems, nor command the states' officers---or those of their political subdivisions---to administer or enforce a federal regulatory program (\href{https://www.law.cornell.edu/supct/html/95-1478.ZO.html}{Cornell LII, Majority Opinion, 521 U.S. at 935}). This anti-commandeering rule is categorical: it applies regardless of whether policymaking discretion is involved, and no case-by-case balancing of federal interests against state burdens is permitted (\href{https://www.law.cornell.edu/supct/html/95-1478.ZO.html}{Cornell LII, Majority Opinion, 521 U.S. at 932, 935}). Accordingly, the Brady Act's interim background-check mandate, 18 U.S.C. \S\ 922(s)(2), and the implicit duty to accept Brady Forms, \S\S\ 922(s)(1)(A)(i)(III)--(IV), are unconstitutional (\href{https://www.law.cornell.edu/supct/html/95-1478.ZS.html}{Cornell LII, Syllabus}).

\textbf{Holding vs. Dicta:} The holding invalidated \S\ 922(s)(2)'s mandatory background-check obligation and the implicit obligation to receive Brady Forms. The Court expressly held that the Act's record-destruction requirement (\S\ 922(s)(6)(B)(i)) and written-reasons requirement (\S\ 922(s)(6)(C)) were not unconstitutional as applied to non-participating CLEOs, but simply inoperative (\href{https://www.law.cornell.edu/supct/html/95-1478.ZS.html}{Cornell LII, Syllabus}). Crucially, the Court declined to decide two issues, leaving them open: (1) whether the remaining obligations on firearms dealers to forward Brady Forms and observe the five-day waiting period were severable, because no dealer or purchaser was a party before the Court; and (2) whether purely ministerial federal reporting requirements imposed on state officials (such as reporting missing children) violate the anti-commandeering rule (\href{https://www.law.cornell.edu/supct/html/95-1478.ZO.html}{Cornell LII, Majority Opinion}; \href{https://library.oconnorinstitute.org/supreme-court/printz-v-united-states-1996/}{O'Connor Library}). Justice Scalia's extended historical critique of Hamilton's nationalist views in \textit{The Federalist} No. 27 and his dismissal of comparative European federalism in footnote 11 reinforced the structural holding but sparked independent doctrinal debate.

\section{Analysis of Opinions}

\subsection{Majority Opinion (Justice Scalia)}
Acknowledging that ``there is no constitutional text speaking to this precise question,'' Justice Scalia organized the Court's analysis around three pillars: historical understanding and practice, constitutional structure, and Supreme Court precedent (\href{https://www.law.cornell.edu/supct/html/95-1478.ZO.html}{Cornell LII, Majority Opinion}).

First, examining \textbf{historical practice}, Scalia concluded that early Congresses never assumed power to commandeer state executive officers. Early statutes assigning duties to state courts (such as naturalization proceedings or maritime seaworthiness determinations) reflected the unique constitutional position of state \textit{judges}, who were expressly bound by the Supremacy Clause (Article VI, Clause 2) and contemplated by the Madisonian Compromise of Article III (\href{https://www.law.cornell.edu/supct/html/95-1478.ZO.html}{Cornell LII, Majority Opinion}). The Extradition Act of 1793 merely implemented the explicit text of the Extradition Clause (Article IV, Section 2) (\href{https://www.law.cornell.edu/supct/html/95-1478.ZO.html}{Cornell LII, Majority Opinion}). Conversely, when the First Congress wanted states to hold federal prisoners in state jails, it issued a \textit{recommendation} to state legislatures rather than a command; when Georgia refused, Congress authorized federal marshals to rent temporary jails (\href{https://www.law.cornell.edu/supct/html/95-1478.ZO.html}{Cornell LII, Majority Opinion}). Turning to \textit{The Federalist}, Scalia read Nos. 36 and 45 as assuming voluntary state consent, and interpreted Hamilton's famous statement in No. 27---that state officers would be ``incorporated'' and ``rendered auxiliary''---to mean only that state officers must not obstruct supreme federal law, corroborated by Madison's narrow examples in No. 44 (\href{https://www.law.cornell.edu/supct/html/95-1478.ZO.html}{Cornell LII, Majority Opinion}).

Second, examining \textbf{constitutional structure}, Scalia anchored the decision in two structural principles. Under \textit{dual sovereignty}, the Framers rejected the Articles of Confederation's model of governing through states and instead empowered Congress to regulate individuals directly, leaving states with ``a residuary and inviolable sovereignty'' that creates Madison's ``double security'' for liberty (\href{https://www.law.cornell.edu/supct/html/95-1478.ZO.html}{Cornell LII, Majority Opinion}). Allowing Congress to impress state police officers into federal service at no cost to the federal treasury would shatter political accountability, letting Congress take credit for solving problems while forcing states to absorb the financial costs and voter blame (\href{https://www.law.cornell.edu/supct/html/95-1478.ZO.html}{Cornell LII, Majority Opinion}). Furthermore, under \textit{Article II separation of powers}, the Take Care Clause vests law execution in a unitary President; transferring execution of federal law to thousands of local CLEOs beyond presidential appointment, removal, or control would shatter executive unity (\href{https://www.law.cornell.edu/supct/html/95-1478.ZO.html}{Cornell LII, Majority Opinion}). Scalia rejected the dissent's reliance on the Necessary and Proper Clause, reasoning that a law violating state sovereignty is not ``proper'' for carrying into execution the Commerce Clause (\href{https://www.law.cornell.edu/supct/html/95-1478.ZO.html}{Cornell LII, Majority Opinion}).

Third, examining \textbf{precedent}, Scalia held that \textit{New York v. United States}, 505 U.S. 144 (1992), controlled categorically. He rejected the Government's distinction between legislative ``policymaking'' and executive ``implementation'' as unworkable and even more intrusive on state autonomy---reducing state officers to ``puppets of a ventriloquist Congress'' (\href{https://www.law.cornell.edu/supct/html/95-1478.ZO.html}{Cornell LII, Majority Opinion}). He likewise rejected the dissent's distinction between commanding states and commanding individual local officers, observing that commanding state officers in their official capacities is indistinguishable from commanding the state itself, and refused to apply any interest-balancing test (\href{https://www.law.cornell.edu/supct/html/95-1478.ZO.html}{Cornell LII, Majority Opinion}).

\subsection{Concurring Opinion (Justice O'Connor)}
Justice O'Connor joined the Court's opinion and wrote briefly to emphasize the ruling's practical boundaries. She noted that the decision did not end the Brady Act's objectives because CLEOs could continue participating voluntarily, the interim provisions expired on November 30, 1998, and Congress remained free to induce state participation contractually through the Spending Clause (\href{https://library.oconnorinstitute.org/supreme-court/printz-v-united-states-1996/}{O'Connor Library}). She also stressed that the Court ``appropriately refrains from deciding whether other purely ministerial reporting requirements imposed by Congress on state and local authorities pursuant to its Commerce Clause powers are similarly invalid,'' citing 42 U.S.C. \S\ 5779(a) (requiring state and local agencies to report missing children to the Department of Justice) (\href{https://library.oconnorinstitute.org/supreme-court/printz-v-united-states-1996/}{O'Connor Library}).

\subsection{Concurring Opinion (Justice Thomas)}
Justice Thomas joined the Court's opinion in full but wrote separately to underscore two additional constitutional limits. First, reiterating his concurrence in \textit{United States v. Lopez}, 514 U.S. 549, 584 (1995), Thomas argued that the Commerce Clause's original meaning does not empower Congress to regulate wholly intrastate, point-of-sale firearm transfers; lacking that underlying power, Congress necessarily lacked authority to conscript state officers to enforce it (\href{https://www.law.cornell.edu/supct/html/95-1478.ZC1.html}{Cornell LII, Thomas Concurrence}). Second, Thomas noted that although the parties had not raised the Second Amendment, growing historical scholarship suggested it protects an individual right to keep and bear arms that could independently bar federal regulation of purely intrastate firearm sales or possession---foreshadowing the Court's later landmark decision in \textit{District of Columbia v. Heller} (\href{https://www.law.cornell.edu/supct/html/95-1478.ZC1.html}{Cornell LII, Thomas Concurrence}).

\subsection{Dissenting Opinion (Justice Stevens)}
Justice Stevens, joined by Justices Souter, Ginsburg, and Breyer, argued that the Constitution's text, history, structure, and precedent all authorized Congress to enlist local officers temporarily to combat an ``epidemic of gun violence'' (\href{https://www.law.cornell.edu/supct/html/95-1478.ZD.html}{Cornell LII, Stevens Dissent}). Textually, the Commerce Clause and Necessary and Proper Clause grant ample affirmative authority, while the Tenth Amendment imposes no independent limit on delegated powers; meanwhile, Article VI's Supremacy and Oath Clauses affirmatively bind state executive and judicial officers to federal law (\href{https://www.law.cornell.edu/supct/html/95-1478.ZD.html}{Cornell LII, Stevens Dissent}). Stevens warned that the majority's categorical rule would cripple Congress during national emergencies such as epidemics, military drafts, or terrorist threats (\href{https://www.law.cornell.edu/supct/html/95-1478.ZD.html}{Cornell LII, Stevens Dissent}).

Historically, Stevens contended that the Framers in \textit{The Federalist} Nos. 27, 36, and 45 explicitly anticipated federal reliance on state officers (such as tax collectors) to avoid building a massive federal bureaucracy, and that early Congresses routinely assigned functionally executive tasks---like naturalization recordkeeping and maritime inspections---to state judges and clerks (\href{https://www.law.cornell.edu/supct/html/95-1478.ZD.html}{Cornell LII, Stevens Dissent}). Structurally, relying on \textit{Garcia v. San Antonio Metropolitan Transit Authority}, 469 U.S. 528 (1985), and Congress's recent Unfunded Mandates Reform Act of 1995, Stevens maintained that the political safeguards of federalism adequately protect state interests (\href{https://www.law.cornell.edu/supct/html/95-1478.ZD.html}{Cornell LII, Stevens Dissent}). He also criticized the majority's Article II theory as contradicting \textit{New York}'s approval of cooperative federalism, and distinguished \textit{New York} as involving coerced state \textit{legislation} rather than modest interim duties on local officers, citing \textit{FERC v. Mississippi}, \textit{Puerto Rico v. Branstad}, and \textit{Testa v. Katt} (\href{https://www.law.cornell.edu/supct/html/95-1478.ZD.html}{Cornell LII, Stevens Dissent}).

\subsection{Dissenting Opinion (Justice Souter)}
Justice Souter joined Stevens's dissent with qualifications, explaining that while early statutory practice was inconclusive, \textit{The Federalist} decisively resolved the case for him (\href{https://www.law.cornell.edu/supct/html/95-1478.ZD1.html}{Cornell LII, Souter Dissent}). Souter found the ``most straightforward reading'' of Hamilton's \textit{Federalist} No. 27---supported by Madison's No. 44 and consistent with Nos. 36 and 45---established that Congress could enlist state executive officers as ``auxiliaries'' to enforce valid federal laws, rejecting Scalia's attempt to reduce ``incorporated'' to mere non-obstruction (\href{https://www.law.cornell.edu/supct/html/95-1478.ZD1.html}{Cornell LII, Souter Dissent}). Souter reconciled this reading with \textit{New York v. United States} (which he had joined) by distinguishing legislative discretion, which cannot be commanded, from executive law enforcement, which can; he added, however, that he would remand to determine whether Congress was obligated to pay fair compensation for state administrative support (\href{https://www.law.cornell.edu/supct/html/95-1478.ZD1.html}{Cornell LII, Souter Dissent}).

\subsection{Dissenting Opinion (Justice Breyer)}
Justice Breyer, joined by Justice Stevens, added a comparative constitutional perspective. Observing that the federal systems of Switzerland, Germany, and the European Union rely on constituent states rather than central bureaucracies to implement central laws, Breyer argued that this empirical experience shows local implementation can actually \textit{preserve} local autonomy and individual liberty better than expanding a national police bureaucracy (\href{https://www.law.cornell.edu/supct/html/95-1478.ZD2.html}{Cornell LII, Breyer Dissent}). Because the Constitution is silent and the Brady Act's ``reasonable effort'' standard prevented overwhelming local agencies, Breyer saw no reason to erect an inflexible categorical bar (\href{https://www.law.cornell.edu/supct/html/95-1478.ZD2.html}{Cornell LII, Breyer Dissent}).

\section{Examples: Future Applications}
\begin{enumerate}
\item \textbf{Same-Side (Invalidated under Anti-Commandeering):} Suppose Congress enacts a federal statute requiring state and municipal police officers to detain undocumented immigrants for 48 hours beyond their scheduled release dates upon receiving a federal immigration detainer request, without conditioning federal funding or providing a voluntary opt-in. Under \textit{Printz}, this direct command to state and local executive officers to administer and enforce a federal immigration program violates the Tenth Amendment's categorical anti-commandeering rule.
\item \textbf{Same-Side (Invalidated under Anti-Commandeering):} Suppose Congress passes a statute prohibiting state legislatures from repealing their existing state-law bans on sports gambling (mirroring \href{https://www.law.cornell.edu/supremecourt/text/16-476}{\textit{Murphy v. NCAA}}, 584 U.S. 453 (2018)). Even though the directive is phrased as a negative prohibition rather than an affirmative command, it directly dictates to state governments how they must exercise their sovereign lawmaking authority rather than regulating private individuals directly, violating the \textit{New York}--\textit{Printz} anti-commandeering doctrine.
\item \textbf{Opposite-Side (Upheld as Valid Spending Condition / Cooperative Federalism):} Suppose Congress offers states \$500 million in federal law-enforcement grants conditioned on state police departments voluntarily conducting five-day background checks on private firearm transfers, leaving states a genuine, non-coercive choice to decline the funds. As Justice O'Connor's concurrence and \textit{New York v. United States} confirm, Congress may use its Spending Clause power to induce voluntary state executive participation without violating \textit{Printz}.
\item \textbf{Opposite-Side (Upheld as Generally Applicable Regulation of State Conduct):} Suppose Congress enacts a federal privacy law restricting both state departments of motor vehicles and private data brokers from selling or disclosing personal driver information without consent (mirroring \textit{Reno v. Condon}, 528 U.S. 141 (2000), discussed in the \href{https://constitution.congress.gov/browse/essay/amdt10-4-2/ALDE_00013627/}{Constitution Annotated}). Because the statute regulates the state's own commercial activity as a database owner rather than requiring state officials to regulate private citizens or enforce a federal program against third parties, it does not violate \textit{Printz}.
\item \textbf{Fence-Sitter (Borderline / Contested):} Suppose Congress passes a standalone statute under the Commerce Clause requiring every local sheriff's office in the nation to compile and transmit monthly statistical reports on missing children or domestic-violence arrests to the U.S. Department of Justice, without attaching federal funds or imposing the same duty on private entities. Justice Scalia's majority opinion expressly reserved the validity of pure information-reporting mandates, and Justice O'Connor's concurrence flagged them as potentially distinct; lower courts and scholars remain divided over whether purely ministerial data-reporting duties cross \textit{Printz}'s bright-line ban on conscripting state executive personnel.
\end{enumerate}

\section{Critique}
\textit{Printz v. United States} is one of the defining landmarks of the Rehnquist Court's ``federalism revolution,'' attracting vigorous praise and criticism from both conservative-originalist and progressive-doctrinal scholars (\href{https://constitution.congress.gov/browse/essay/amdt10-4-2/ALDE_00013627/}{Constitution Annotated}).

\textbf{Originalist and Textualist Perspectives:} Formalist and structural-originalist defenders praise \textit{Printz} for enforcing the Constitution's foundational architecture of dual sovereignty and preventing Congress from eroding state independence through unfunded administrative conscription (\href{https://brooklynworks.brooklaw.edu/context/faculty/article/2006/viewcontent/Gold22.pdf}{Gold, \textit{Formalism and State Sovereignty in Printz v. United States}}). Yet \textit{Printz} has also drawn sharp criticism \textit{on originalist and textualist grounds}. Scholars note the irony that Justice Scalia---the Court's foremost textualist---explicitly admitted ``there is no constitutional text speaking to this precise question'' and relied heavily on structural inference and historical silence (\href{https://journals.library.wustl.edu/lawreview/article/5537/galley/22370/view/}{Washington Univ. Law Review, \textit{The Constitution According to Justices Scalia and Thomas}}). Critics such as Gene Nichol and Saikrishna Prakash have questioned Scalia's dismissal of Hamilton's \textit{Federalist} No. 27 and early congressional practice, arguing that the historical record was at best ambiguous and arguably favored federal authority to enlist state executive officers---making Scalia's categorical prohibition a departure from strict textualism (\href{https://scholarlycommons.law.case.edu/cgi/viewcontent.cgi?article=1253&context=caselrev}{Case Western Reserve L. Rev., citing Nichol (1999)}; \href{https://scholarship.law.wm.edu/cgi/viewcontent.cgi?article=1572&context=wmborj}{Wm. \& Mary Bill Rts. J., citing Prakash (1993)}).

\textbf{Progressive and Doctrinal Perspectives:} Progressive and functionalist scholars, such as Vicki C. Jackson and Evan Caminker, critique \textit{Printz}'s rigid, anti-functionalist bright line as poorly grounded in history and counterproductive to workable federalism (\href{https://brooklynworks.brooklaw.edu/context/faculty/article/2006/viewcontent/Gold22.pdf}{BrooklynWorks, citing Jackson, 111 Harv. L. Rev. 2180 (1998), and Caminker, 95 Colum. L. Rev. 1001 (1995)}). Echoing Justices Stevens and Breyer, functionalists argue that forbidding modest federal enlistment of state officials incentivizes Congress to expand centralized federal bureaucracies, ironically diminishing local control, and that the Court's ``political accountability'' rationale underestimates voters' ability to identify federal mandates---especially when vocal local officials like Sheriffs Printz and Mack publicly denounce them. At the same time, in the decades since \textit{Printz}, progressive jurisdictions have embraced the anti-commandeering doctrine as a vital shield protecting ``sanctuary'' states and cities from federal immigration enforcement mandates, demonstrating the doctrine's cross-ideological durability (\href{https://texaslawreview.org/making-federalism-great-again/}{Texas Law Review, \textit{Making Federalism Great Again}}).

\textbf{Competing Underlying Values:} At its core, \textit{Printz} pits \textit{formal structural federalism, political accountability, and unitary executive control} against \textit{pragmatic cooperative federalism, national regulatory efficiency, and democratic deference to Congress}. By choosing a bright-line categorical prohibition over a balancing test, the Court prioritized clear constitutional boundaries between sovereigns even at the cost of limiting flexible intergovernmental responses to national problems.

\section{Key Quotations}

\begin{quote}
``We held in New York that Congress cannot compel the States to enact or enforce a federal regulatory program. Today we hold that Congress cannot circumvent that prohibition by conscripting the State's officers directly. The Federal Government may neither issue directives requiring the States to address particular problems, nor command the States' officers, or those of their political subdivisions, to administer or enforce a federal regulatory program. It matters not whether policymaking is involved, and no case by case weighing of the burdens or benefits is necessary; such commands are fundamentally incompatible with our constitutional system of dual sovereignty.''
\end{quote}
\noindent --- Justice Antonin Scalia, Opinion of the Court, 521 U.S. at 935 (\href{https://www.law.cornell.edu/supct/html/95-1478.ZO.html}{Cornell LII}; \href{https://www.courtlistener.com/opinion/118148/printz-v-united-states/}{CourtListener}).

\begin{quote}
``The power of the Federal Government would be augmented immeasurably if it were able to impress into its service--and at no cost to itself--the police officers of the 50 States.''
\end{quote}
\noindent --- Justice Antonin Scalia, Opinion of the Court, 521 U.S. at 922 (\href{https://www.law.cornell.edu/supct/html/95-1478.ZO.html}{Cornell LII}; \href{https://www.courtlistener.com/opinion/118148/printz-v-united-states/}{CourtListener}).

\begin{quote}
``When a `La[w] . . . for carrying into Execution' the Commerce Clause violates the principle of state sovereignty reflected in the various constitutional provisions we mentioned earlier, supra, at 919, it is not a `La[w] . . . proper for carrying into Execution the Commerce Clause,' and is thus, in the words of The Federalist, `merely [an] ac[t] of usurpation' which `deserve[s] to be treated as such.'''
\end{quote}
\noindent --- Justice Antonin Scalia, Opinion of the Court, 521 U.S. at 923--924 (\href{https://www.law.cornell.edu/supct/html/95-1478.ZO.html}{Cornell LII}; \href{https://www.courtlistener.com/opinion/118148/printz-v-united-states/}{CourtListener}).

\begin{quote}
``There is not a clause, sentence, or paragraph in the entire text of the Constitution of the United States that supports the proposition that a local police officer can ignore a command contained in a statute enacted by Congress pursuant to an express delegation of power enumerated in Article I.''
\end{quote}
\noindent --- Justice John Paul Stevens, Dissenting, 521 U.S. at 944 (\href{https://www.law.cornell.edu/supct/html/95-1478.ZD.html}{Cornell LII}; \href{https://www.courtlistener.com/opinion/118148/printz-v-united-states/}{CourtListener}).

\begin{quote}
``At least some other countries, facing the same basic problem, have found that local control is better maintained through application of a principle that is the direct opposite of the principle the majority derives from the silence of our Constitution. The federal systems of Switzerland, Germany, and the European Union, for example, all provide that constituent states, not federal bureaucracies, will themselves implement many of the laws, rules, regulations, or decrees enacted by the central `federal' body.''
\end{quote}
\noindent --- Justice Stephen G. Breyer, Dissenting, 521 U.S. at 976 (\href{https://www.law.cornell.edu/supct/html/95-1478.ZD2.html}{Cornell LII}; \href{https://www.courtlistener.com/opinion/118148/printz-v-united-states/}{CourtListener}).

\end{document}
